%!TEX root = ../../main.tex
\section{사건의 가치를 승률 변화로 매기기}
\label{sec:rel-eventvalue}

앞 절의 교환가치는 킬과 오브젝트에 손으로 정한 가중을 곱한 점수였다. 가중을 손으로 정하지 않는 길이 있다. 게임 안의 사건 하나가 승률을 올린 만큼의 가치를 갖는다고 보는 것이며, 이 생각은 리그 오브 레전드 연구에 이미 있다. Maymin \cite{maymin2021smart}은 로지스틱 승률 모형으로 킬과 데스를 ``영리한 킬''과 ``무가치한 데스''로 나누었다. 축구에서 선수의 행동 하나하나를 득점 확률의 변화로 평가하는 연구 \cite{DecroosEtAl2019}도 같은 잣대를 쓴다.

변화를 다룰 때에는 방향과, 부호를 살린 평균과, 변화의 크기를 서로 구분할 필요가 있다. 이 구분은 금융의 사건 연구에 있다. 그 연구는 사건 전후의 수익을 평균 비정상수익으로 검정하는 한편, 사건이 수익의 변동 자체를 바꾸는 경우를 따로 다룬다 \cite{BrownWarner1985}. 사건이 어떤 종목의 수익은 올리고 어떤 종목의 수익은 내리면 평균은 0에 가까우면서 변동만 커지므로, 평균을 전제로 한 검정은 그 반응을 놓친다 \cite{BoehmerEtAl1991}. 방향은 또 다른 양이다. 변화 방향을 예측하는 연구는 평균은 예측할 수 없어도 방향은 예측할 수 있는 경우가 있고 그 반대의 경우도 있음을 보였다 \cite{ChristoffersenDiebold2006}.

본 논문은 이 잣대를 사건 하나에서 교전 한 구간으로 옮긴다. 그리고 같은 구분을 따라 교전 구간의 추정 승률 변화를 세 양으로 나눈다. 부호를 살린 평균, 변화의 크기, 그리고 변화가 오르는 쪽인지의 방향이다(제 \ref{sec:value-triad} 절). 이 가운데 본 논문이 예측 목표로 삼는 것은 방향이며, 교전이 시작될 때의 승률과 경기 시각만으로 알 수 있는 부분은 기준선에 맡기고 그것을 넘어서는 부분만 센다. Maymin의 연구는 이 게임에서 승률 변화를 다룬 직접적 선행 연구이며, 본 논문은 세 가지가 다르다. 단위가 교전 구간이고, 승률 모형을 학습한 뒤 고정하며, 승률과 시간만 쓰는 기준선과 견준다. 금융의 사건 연구는 세 양을 나누어 본다는 방법상의 유사성 때문에 인용한다.
